\section{CLAUDE.md 与 AGENTS.md：Agent 配置文件}

\subsection{基本概念}

在触碰任何其他 harness 配置点之前，通常值得先定制你的 \texttt{CLAUDE.md} / \texttt{AGENTS.md} 文件。这些是放在仓库根目录的 markdown 文件，harness 会确定性地将其注入 agent 的 system prompt。

\subsection{ETH Zurich 研究的启示}

ETH Zurich 发表了一项研究，测试了 138 个 agentfile 在多个仓库中的效果，结论是大多数 agentfile ``无用甚至有害''。研究的主要发现：

\begin{table}[H]
\centering
\caption{ETH Zurich Agentfile 研究关键发现}
\begin{tabular}{@{}p{5cm}p{8cm}@{}}
\toprule
\textbf{发现} & \textbf{具体数据} \\
\midrule
LLM 生成的文件 & 实际损害表现，且成本增加 20\%+ \\
人工编写的文件 & 仅帮助约 4\% \\
推理 token 消耗 & Agent 多花 14--22\% 的推理 token 处理指令 \\
步骤与工具调用 & 更多步骤、更多工具调用，但分辨率未提升 \\
代码库概览 & 目录列表完全无帮助，agent 自己就能发现仓库结构 \\
\bottomrule
\end{tabular}
\end{table}

\subsection{编写高质量 Agentfile 的原则}

HumanLayer 认为该研究恰好验证了他们之前的建议：

\begin{importantbox}{Agentfile 最佳实践}
\begin{enumerate}[leftmargin=1.5em]
    \item \textbf{避免自动生成}：LLM 生成的 agentfile 效果更差
    \item \textbf{指令越少越好}：不省略必要指令的前提下，尽量精简
    \item \textbf{使用渐进式披露}（Progressive Disclosure）：不要把所有信息都塞进去
    \item \textbf{保持简洁和普适}：内容应普遍适用，避免过多条件规则
    \item \textbf{不要过度引导工具使用}：过度指定``用什么工具''反而导致更差的结果
\end{enumerate}
HumanLayer 自己的 CLAUDE.md 不到 60 行。
\end{importantbox}

\begin{warningbox}{常见误区}
\begin{itemize}[leftmargin=1.5em]
    \item 把代码库概览和目录结构放进 agentfile --- agent 自己会探索
    \item 用 LLM 自动生成 agentfile --- 实际效果为负
    \item 塞入大量条件规则（``如果 X 则 Y，否则 Z''）--- 增加推理负担却不提高成功率
    \item 详细指定每个任务应该使用什么工具 --- 过度约束模型的灵活性
\end{itemize}
\end{warningbox}

\subsection{本章小结}

CLAUDE.md / AGENTS.md 是最基础的 harness 配置点。但``有配置''不等于``好配置''：自动生成的、信息过载的、过度引导的文件反而有害。核心原则是精简、普适、渐进式披露。


